\documentclass[11pt]{article}

\usepackage[margin=0.8in]{geometry}
\usepackage[T1]{fontenc}
\usepackage[utf8]{inputenc}
\usepackage{lmodern}
\usepackage{enumitem}
\usepackage{hyperref}
\usepackage{titlesec}

\setlength{\parindent}{0pt}
\setlength{\parskip}{0.4em}

\titleformat{\section}
  {\Large\bfseries}
  {}
  {0pt}
  {}

\titlespacing{\section}
  {0pt}
  {1.5em}
  {0.5em}

\newcommand{\sectionline}{\noindent\rule{\linewidth}{0.5pt}}

\begin{document}

\begin{center}
    {\LARGE \textbf{Jegannarayanan Sivakasiulaganathan}} \\
    Dedham, MA \quad | \quad +1-669-233-1124 \quad | \quad Jegannarayanan\_Sivakasiulaganathan@student.uml.edu
\end{center}

Graduate Computer Science student focused on compilers and programming 
languages, building compilers in my personal time and working with testing and 
verification in academic settings.

Expecting to graduate in May 2027, focused on compilers, programming languages, 
and deep learning systems. Software developer with 7 years of experience 
across desktop and full stack software engineering. 

\section*{Education}
\sectionline

\textbf{University of Massachusetts Lowell} \\
M.S. Computer Science \hfill Expected May 2027

% TODO: Put the what and why in some of these
\textbf{Relevant Graduate Coursework}
\begin{itemize}[leftmargin=*]
    \item COMP 7030 CS Research -- Rocq (currently taking)
    \item COMP 5300 Effective Functional Programming (currently taking)
    \item COMP 5300 Deep Learning: LLMs and LVMs (currently taking)
    \item COMP 5030 Graduate Algorithms
    \item COMP 5150 Operating Systems I
    \item COMP 5430 Artificial Intelligence
\end{itemize}

\textbf{University of Massachusetts Amherst} \\
B.S. Computational and Applied Mathematics \hfill May 2018

\textbf{Relevant Undergraduate Coursework}
\begin{itemize}[leftmargin=*]
    \item \textbf{Theory \& Foundations:} Computation, Algorithms, Formal Language Theory, Logic in Computer Science
    \item \textbf{Mathematics:} Multivariable Calculus, Linear Algebra, Ordinary Differential Equations, Abstract Algebra, Advanced Calculus
    \item \textbf{Applied Math \& Scientific Computing:} Applied Linear Algebra, Scientific Computing, Mathematical Modeling
\end{itemize}

\section*{Technical Skills}
\sectionline

\textbf{General:} Compilers, Deep Learning, Formal Verification \\
\textbf{Languages:} C++, Python, Rocq, Haskell, C\#, Typescript \\
\textbf{Systems:} Linux \\
\textbf{Tools:} Docker, GitLab CI/CD, Git \\
\textbf{Scripting:} Bash, Python \\
\textbf{GitHub:} \url{https://github.com/handel-fan}

\section*{Personal Projects}
\sectionline

\begin{itemize}[leftmargin=*]
    \item Porting an interpreter for the Monkey programming language from C++ to Go \\
    \url{https://github.com/handel-fan/first-interpreter-in-cpp}
    \item Built LLVM from source and explored compiler toolchains
    \item Benchmarked Python matrix multiplication techniques and performance tradeoffs
\end{itemize}

\section*{Work Experience}
\sectionline

% TODO: Add impact here
\textbf{KBRA} \hfill Dresher, PA (Remote) \\
\textit{Software Developer} \hfill March 2023 -- Present
\begin{itemize}[leftmargin=*]
    \item Led technical debt initiatives to improve code quality and maintainability
    \item Contributed to full stack development, including React frontend work
\end{itemize}

\textit{Associate Software Developer} \hfill August 2022 -- March 2023
\begin{itemize}[leftmargin=*]
    \item Developed backend services and APIs
    % TODO: Which ones?
    \item Worked across multiple technologies in a production environment
\end{itemize}

\textbf{Aspen Technology} \hfill Bedford, MA \\
\textit{Software Developer} \hfill July 2021 -- August 2022
\begin{itemize}[leftmargin=*]
    \item Contributed to desktop software development and maintenance
    \item Resolved bugs and supported product stability
\end{itemize}

\textit{Software Developer I} \hfill May 2018 -- July 2021
\begin{itemize}[leftmargin=*]
    \item Provided support for desktop applications and implemented bug fixes
\end{itemize}

\end{document}
